\section{结论}

\subsection{四问主要结果}

\begin{table}[htbp]
	\caption{四问关键指标汇总（本文结果）}
	\label{tab:allmetrics}
	\centering
	\zihao{-5}
	\begin{tabular}{p{1.5cm}p{4.0cm}p{5.0cm}p{3.6cm}}
		\toprule
		问题 & 模型 & 主要结果 & 证据等级 \\
		\midrule
		问题一
		& 地形约束非线性最大安全载荷 $+$ 质量—体积二维装箱 $+$ Martello--Toth 解析下界
		& \textbf{18 架次}（全 C 型）/ \SI{75.07}{kWh} / 累计 \SI{9.37}{h}；逐区等于下界
		& \textbf{架次数可证最优} \\
		问题二
		& 异构多架次带时间窗 VRPTW $+$ 资源约束调度
		& \textbf{25 架次}（A8$+$B10$+$C7，8 架全投入）/ \SI{77.30}{kWh} / 完工 \SI{3.02}{h} / 准时率 \textbf{100.0\%}
		& 经验证的可行解（距下界 24 差 1） \\
		问题三
		& 连续通信约束 $+$ 轨迹逐时刻采样 $+$ 悬停点选址 $+$ 窗口硬约束排班
		& \textbf{时间轴覆盖 $11/20=55.0\%$}（几何可达 $20/20$）/ 中继 6 架次 / \SI{79.65}{kWh}
		& 硬违规 0 条；\textbf{中继缺口 4 架} \\
		问题四
		& 图连通分量（原子单元）$+$ 组内资源峰值核算
		& $K\in[1,3]$，\textbf{$K=2$/$K=3$ 均可零架次改动}；资源 $33/34/38$
		& \textbf{精确图论结论} \\
		\bottomrule
	\end{tabular}
\end{table}

\subsection{主要结论}

\textbf{（1）结构载重是问题一的主导约束，能量只是次要约束。}A 型在全部 15 个服务区都受结构上限 $\Qg$ 约束，B 型 14/15，仅 C 型在 5 个远距离服务区受能量约束。因此单点往返场景下的组批方案主要由最大载货质量与装载体积决定。附件取 $\rhog=0.20$ 恰在效率最优区间；$\rhog\ge0.375$ 时远距离服务区即使空载也无法返航。

\textbf{（2）首批保障时限可以完全满足，其可行性由“供给—需求对偶”决定。}60\,min 档服务区恰为 8 个，与题目机队 8 架一一对应；只要把首批专架次按机队槽位轮转铺开、并在派发时以最小松弛优先抢占首轮机位，8 个区即可同时开工。实测最晚首批交付 \SI{3587}{s}（裕度 \SI{12}{s}），准时率 \textbf{100.0\%}。\textbf{“不可行”的旧结论源于机队坍缩与派发退化两处建模缺陷。}

\textbf{（3）中继无人机是必需项，但题目配置的中继数量不足。}25 个运输架次中 20 个存在直连中断（平均占比 28.4\%），故中继不可省略。然而单架中继的在站时长上限仅约 \SI{140}{min}，而 20 个架次的真实中断时长合计约 \SI{168.5}{min}、最大并发 6 架次，2 架中继在时间轴上\textbf{无法}保障全部架次。本文按窗口硬约束重排后得到 \textbf{55.0\%} 的时间轴覆盖率，并如实报告 \textbf{4 架中继缺口}。

\textbf{（4）覆盖率必须用时间轴口径。}几何可达覆盖率（$20/20$）只说明“存在合适的悬停点”，是选址能力的\textbf{上界}；时间轴覆盖率（$11/20$）才回答“中继那一刻是否在站”。\textbf{不得把前者当作“已保障”上报。}

\textbf{（5）任务分区不省资源，但能降低单组峰值工作量。}合法组数受原子单元拓扑支配（$K\in[1,3]$），$K=2$ 与 $K=3$ 均可零架次改动实现。资源总量随 $K$ 上升（$33\to34\to38$），组间不均衡度上升，但单组最大工作量由 \SI{13.95}{h} 降至 \SI{13.00}{h}。分区的价值在于\textbf{削峰}而非省资源。

\subsection{资源缺口与建议}

本题最关键的资源短板是\textbf{中继无人机}：库存 2 架，而问题三需 6 个中继架次（缺口 \textbf{4 架}）；该缺口在三种分区下都存在且分区无法缓解。其次是共享电池（整队一组时 A/B/C 三型各缺约 1 组）。

\textbf{建议}：\textbf{①}把中继无人机增至 4$\sim$6 架，可直接消除问题三的覆盖瓶颈；\textbf{②}增加共享电池组以缓解周转压力；\textbf{③}在资源受限时，用时空联合优化把有限的在站时长优先投给“单位时间可覆盖架次数最多”的悬停点。

\subsection{全文证据状态声明}

为保证表述的严谨性，本文明确声明：表~\ref{tab:allmetrics} 中的全部数值均为\textbf{本文程序求解并经独立校验器复核的结果}；其中\textbf{问题一的 18 架次}与\textbf{问题四的 $K$ 取值范围}属可证最优/精确结论，\textbf{问题二的 25 架次}仅为经验证的可行解、\textbf{不声称全局最优}，\textbf{问题三的 55.0\%} 覆盖率是\textbf{题目资源约束下的可达上界}而非算法能力上限。附件直接核对的输入参数已在 \S\ref{sec:data} 与表~\ref{tab:uav} 列出，未引入任何外部数据集。
